\documentclass[11pt]{article}
\usepackage[a4paper,margin=21mm]{geometry}
\usepackage{amsmath,amssymb,amsthm,hyperref}
\newtheorem{theorem}{Theorem}
\newtheorem{corollary}[theorem]{Corollary}
\newtheorem{proposition}[theorem]{Proposition}
\title{GST emergent complex geometry and causal topology: a dual realization of Hodge alignment and mixed-weight birth}
\date{23 October 2026}
\begin{document}\maketitle
\begin{abstract}
We test the owner's topology-first, dimension-emergent spatial program,
rather than beginning with an algebraic Chow-surjectivity hypothesis.
One arbitrary finite GST causal graph is given two \emph{actual}
algebro-geometric realizations over \emph{every} original smooth projective
complex $X$: (i) a family of smooth projective products with variable
elliptic complex structures and entirely pure Hodge cohomology, and
(ii) a singular nodal graph-gluing realization with topology-born
weight-zero $H^1$. We derive exact complementary laws: the complex
rank of Hodge-alignment constraints on rational topological edge
classes equals $|V|-c$, whereas the topological $W_0$ birth
dimension equals $|E|-|V|+c$. Their sum is the number of causal edges.
An integral edge class is Hodge precisely when the corresponding
elliptic parameters match; when it is Hodge it is explicitly the
canonical Betti class of a \emph{genuine geometric diagonal divisor},
not of a stipulated Chow selector. The result gives a strong,
non-circular geometry-to-algebraicity theorem for the entire
edge family, and shows why pure-Hodge persistence of the full
smooth projective world and Hodge preservation of a fixed rational
topological class are different mathematical events. No universal
rational Hodge conjecture assumption is made or concluded.
\end{abstract}

\section{Pregeometry and the two typed geometric realizations}
Let $\Gamma=(V,E,s,t)$ be an arbitrary finite loop-free graph.
Parallel edges are allowed. Its underlying unoriented graph
has $c$ connected components. Orientation is only bookkeeping.
The GST carrier of \texttt{GSTGeneralSpace.lean} has no intrinsic
dimension or metric. Its optional realization may separately
select a scheme and a causality-to-geometry map; the
$k$-cell index of \texttt{GSTGraphV2UnboundedHigherCausality.lean}
must not be mistaken for the scheme's complex dimension.

Fix \emph{arbitrary} nonempty smooth projective $X$ over $\mathbb C$.
For each vertex $v$ choose a marked elliptic complex modulus
$\tau_v\in\mathfrak H=\{z\in\mathbb C:\operatorname{Im}z>0\}$ and
$E_{\tau_v}=\mathbb C/(\mathbb Z+\tau_v\mathbb Z)$.
The \emph{smooth whole-world realization} is
\[
 Y_{\boldsymbol\tau}=X\times\prod_{v\in V}E_{\tau_v}.
 \tag{1}
\]
Every $Y_{\boldsymbol\tau}$ is smooth projective, of complex
dimension $\dim_{\mathbb C}X+|V|$, with a pure Hodge structure.
Even when the actual smooth topology is held fixed as
$X\times(S^1)^{2|V|}$, its realized holomorphic Hodge subspaces
vary with $\boldsymbol\tau$.

The \emph{topology-folded realization} is
\[
 T_\Gamma=X\times C_\Gamma,                             \tag{2}
\]
where $C_\Gamma$ is the genuine projective nodal curve obtained
by taking one elliptic curve at each vertex and identifying
distinct smooth points pairwise for every graph edge.
Points for distinct edge ends are chosen disjointly.
The dual graph of $C_\Gamma$ is $\Gamma$.
If $\Gamma$ has an edge then $T_\Gamma$ is singular, not a
new smooth-projective counterexample. Both (1) and (2)
share the original geometric $X$ as an external
component, and use the \emph{same} GST causal incidence graph;
they are different optional realization functors.

\section{Exact Hodge-class break inside pure smooth-projective geometry}
Use real coordinates $x_v,y_v\in\mathbb R/\mathbb Z$
and the marked holomorphic forms
$dz_v=dx_v+\tau_v\,dy_v$.
For an edge $e:v\to w$, pull back the integral degree-two class
\[
 \boxed{
 \Omega_e=(dx_v-dx_w)\wedge(dy_v-dy_w)
      \in H^2(Y_{\boldsymbol\tau},\mathbb Z).
 }                                                            \tag{3}
\]
The topological marking identifies this class across
different smooth-projective $\boldsymbol\tau$.
Only the two elliptic factors at $v,w$ participate.

\begin{theorem}[GST geometric Hodge-alignment/break equation]
For every arbitrary $X,\Gamma,\boldsymbol\tau$, every edge
$e:v\to w$ has the \emph{exact} holomorphic period
\[
 \boxed{
 \Omega_e\wedge dz_v\wedge dz_w
   =(\tau_w-\tau_v)\,
     dx_v\wedge dy_v\wedge dx_w\wedge dy_w
 }                                                            \tag{4}
\]
on its two-factor real four-torus, with the orientation
shown. Consequently,
\[
 \boxed{
 \Omega_e\in H^{1,1}(Y_{\boldsymbol\tau})
 \ \Longleftrightarrow\ \tau_v=\tau_w.
 }                                                            \tag{5}
\]
Whenever (5) holds, the two elliptic factors coincide
as marked complex elliptic curves, and
\[
 \boxed{
 \Omega_e=\mathrm{cl}^{can}_{Y_{\boldsymbol\tau}}
 \bigl(\mathrm{pr}_{v,w}^*[\Delta_{v,w}]\bigr),
 }                                                            \tag{6}
\]
where $\Delta_{v,w}\subset E_{\tau_v}\times E_{\tau_w}$
is the \emph{actual algebraic diagonal divisor}.
Thus the fixed integral topological edge class is
Hodge if and only if its concrete algebraic diagonal
representative exists in this marked geometric realization.
These are implications proved by explicit geometry,
not by assuming the Hodge conjecture.
\end{theorem}
\begin{proof}
Expand
$(dx_v-dx_w)\wedge(dy_v-dy_w)\wedge
 (dx_v+\tau_vdy_v)\wedge(dx_w+\tau_wdy_w)$.
The determinant of the coefficient matrix with rows
$(1,0,-1,0)$, $(0,1,0,-1)$,
$(1,\tau_v,0,0)$, $(0,0,1,\tau_w)$
is exactly $\tau_w-\tau_v$, proving (4).
The real rational degree-two class is of type $(1,1)$
precisely when its $(0,2)$ component (equivalently
the wedge against $dz_v\wedge dz_w$) vanishes:
the two elliptic factors have complex dimension two
and conjugation gives the other off-diagonal component.
Pullback along the projection to those two factors
and restriction along a chosen point of $X$ and all
other elliptic factors show the equivalence in the
full $Y_{\boldsymbol\tau}$.
When $\tau_v=\tau_w$, the diagonal in the two
complex tori is a complex subgroup and hence
an algebraic divisor in their projective product.
Its normal-circle topological Thom class is
exactly $(dx_v-dx_w)\wedge(dy_v-dy_w)$ in the
stated marking. This proves (6), including
the \emph{canonical} algebraic cycle-class comparison.
\end{proof}
\begin{remark}
Even when $\tau_v\ne\tau_w$ the \emph{whole}
$H^\bullet(Y_{\boldsymbol\tau},\mathbb Q)$ remains
a pure Hodge structure. What breaks is the Hodge
status of the \emph{marked rational edge class}.
This is an actual family of smooth projective
geometries, not an abstract rational countermodel.
It neither disproves the rational Hodge conjecture
nor proves it on all $X$.
\end{remark}

\section{Causal-graph law for the full edge Hodge-alignment locus}
Let $B_\Gamma:\mathbb C^V\to\mathbb C^E$ be the
\emph{realized causal graph incidence operator},
\[
 (B_\Gamma\boldsymbol\tau)_e
       =\tau_{t(e)}-\tau_{s(e)}.                          \tag{7}
\]
Unlike the GST dimensionless carrier itself,
$\boldsymbol\tau$ has the complex-geometric meaning
of actual elliptic moduli.

\begin{theorem}[Exact complex-geometric alignment rigidity]
The simultaneous geometric Hodge condition on every
edge class $\Omega_e$ is \emph{exactly}
\[
 \boxed{
  B_\Gamma\boldsymbol\tau=0
   \ \Longleftrightarrow\
  \tau_v=\tau_w
   \text{ whenever }v,w\text{ lie in the same connected
    component of }\Gamma.
 }                                                        \tag{8}
\]
In particular the actual Hodge-alignment locus in
$\mathfrak H^V$ is
$\mathfrak H^{c}$ (diagonally embedded by component),
and the complex codimension of that locus is
\[
 \boxed{\operatorname{rank}_{\mathbb C}B_\Gamma=|V|-c.}  \tag{9}
\]
Every class $\Omega_e$ simultaneously lying in the
rational $(1,1)$ fiber has the explicitly algebraic
diagonal representative (6).
\end{theorem}
\begin{proof}
Equation (5) identifies each complex Hodge condition
with the corresponding edge difference zero.
Such differences vanish on every edge precisely
when $\boldsymbol\tau$ is constant along every graph
path. It can then vary freely, once per connected
component, through the upper half-plane.
Hence $\ker B_\Gamma\cong\mathbb C^c$
and rank-nullity proves (9).
The true diagonal construction proves the last
sentence.
\end{proof}

\section{Genuine topology-folded realization: purity break and duality}
\begin{theorem}[Exact causal graph mixed-weight birth]
Let $\nu:\coprod_v E_{\tau_v}\to C_\Gamma$ be the
normalization of the actual nodal curve in (2).
Then the MHS exact sequence
\[
 0\longrightarrow
 \operatorname{coker}(B_\Gamma:\mathbb Q^V\to\mathbb Q^E)
 \longrightarrow H^1(C_\Gamma,\mathbb Q)
 \longrightarrow\bigoplus_v H^1(E_{\tau_v},\mathbb Q)
 \longrightarrow0                                             \tag{10}
\]
has weight-zero subspace
\[
 \boxed{
 \dim_{\mathbb Q}W_0H^1(T_\Gamma,\mathbb Q)
 =|E|-|V|+c. }                                               \tag{11}
\]
For the product with arbitrary smooth projective $X$,
$H^1(T_\Gamma)=H^1(X)\oplus H^1(C_\Gamma)$;
the former is pure weight one, so it contributes no
weight-zero classes. Hence the true topology fold
breaks pure weight-one Hodge structure precisely
when $\Gamma$ has a graph cycle. If $\Gamma$ is a
forest, this specific topology fold does not birth
weight-zero $H^1$.
\end{theorem}
\begin{proof}
The constant-sheaf normalization exact sequence is
$0\to\mathbb Q_{C_\Gamma}
 \to\nu_*\mathbb Q_{\coprod E_v}
 \to\oplus_{e}\mathbb Q_{q_e}\to0$.
In its long cohomology sequence, the $H^0$-connecting
map is the incidence operator
$B_\Gamma:\mathbb Q^V\to\mathbb Q^E$.
The skyscraper node terms are pure of weight zero;
each smooth projective elliptic $H^1(E_v)$ is
pure of weight one. Exactness and strictness of
the MHS weight filtration identify
$W_0H^1(C_\Gamma)$ with $\operatorname{coker}B_\Gamma$.
Its rank is $|E|-(|V|-c)$.
The product formula for $H^1$ and the purity of
$H^1(X)$ give (11).
\end{proof}

\begin{corollary}[GST causality--Hodge alignment complementarity]
For every finite original GST graph $\Gamma$,
the two different actual geometric realization
experiments satisfy the \emph{exact integer identity}
\[
 \boxed{
 \underbrace{\operatorname{codim}_{\mathbb C}
   \{\text{all smooth edge classes }\Omega_e
     \text{ remain Hodge}\}}_{|V|-c}
 +
 \underbrace{\dim_{\mathbb Q}W_0H^1(T_\Gamma,\mathbb Q)}
   _{|E|-|V|+c}
 =\boxed{|E|}.                                                \tag{12}
\]
This is an identity of ranks of the \emph{same}
causal incidence operator, realized in two genuinely
different geometric categories. It is not an
identification of the moduli conormal space with
the mixed-Hodge weight-zero space, nor does it
postulate a Hodge-to-Chow isomorphism.
\end{corollary}
\begin{proof}
The first term is $\operatorname{rank}B_\Gamma$,
the second is $\operatorname{dim}\operatorname{coker}B_\Gamma$.
Rank-nullity at the edge-coordinate target
$\mathbb Q^E$ proves (12), with matching
integer ranks over $\mathbb Q$ and $\mathbb C$.
\end{proof}


\section{Exclusive GST causal birth-or-break law for each new edge}
\begin{theorem}[One-edge causal dimension/topology bifurcation]
Let $\Gamma$ be \emph{any} finite causal graph
with a fixed vertex set $V$, and insert an edge $e$
between vertices $v,w$ which is not yet present
as that edge (parallel edges are allowed).
Let $a_\Gamma=|V|-c(\Gamma)$ denote the codimension
of the simultaneous whole-edge Hodge-alignment
locus of the \emph{smooth} realization, and
let $b_\Gamma=|E|-|V|+c(\Gamma)$ denote the
weight-zero $H^1$ rank in the \emph{nodal topology fold}.
Then the new geometric world satisfies the sharp
mutually exclusive law
\[
 \boxed{
 (\Delta a,\Delta b)=
 \begin{cases}
 (1,0), & v,w\text{ were in different causal components},\\
 (0,1), & v,w\text{ were already causally connected}.
 \end{cases}
 }                                                              \tag{13}
\]
In particular,
\[
 \boxed{\Delta a+\Delta b=1,\qquad
       (\Delta a)(\Delta b)=0.}                                 \tag{14}
\]
This is a precise algebraic-geometric versus
topology-born mixed-Hodge \emph{bifurcation}:
\begin{itemize}
\item A connecting edge adds one independent
  holomorphic modulus-alignment equation
  $\tau_v=\tau_w$, but does not add a new
  graph-loop weight-zero mode.
\item A cycle-closing edge adds no independent
  smooth-world Hodge-alignment equation, but
  creates exactly one additional
  weight-zero $H^1$ direction in the nodal world.
\end{itemize}
Both statements work for arbitrary finite causality
depth, arbitrary number of components, arbitrary
original smooth projective $X$, and arbitrary
dimension of $X$.
\end{theorem}
\begin{proof}
If $v,w$ lie in different connected components,
the edge merges exactly two components, so
$c(\Gamma+e)=c(\Gamma)-1$.
Therefore
$a_{\Gamma+e}-a_\Gamma=
 (|V|-(c-1))-(|V|-c)=1$
and
$b_{\Gamma+e}-b_\Gamma=
 ((|E|+1)-|V|+(c-1))-(|E|-|V|+c)=0$.
If they were already connected, the new edge
does not change $c$. The new equality
$\tau_v=\tau_w$ follows from equalities along
an existing path and is redundant, so
$\Delta a=0$.
The new edge independently closes one
causal loop, and the normalized nodal curve's
weight-zero first-cohomology rank rises by one:
$\Delta b=1$. This proves (13)--(14)
without invoking any rational Hodge conjecture
assumption or constructing arbitrary Chow preimages.
\end{proof}

\begin{corollary}[Topology graph path-independent net causal law]
For any finite edge insertion sequence from the
edgeless graph on a fixed vertex set, at \emph{every}
intermediate world the total number of
connecting-edge alignment births is exactly
$|V|-c(\Gamma)$, while the number of
cycle-closing topology weight births is
$|E|-|V|+c(\Gamma)$, independent of the
chosen edge order. Reordering causal births
may change \emph{which} edge produces a given
kind of birth, but not either net count.
\end{corollary}
\begin{proof}
Sum the two cases of (13) along the path;
the initial counts both vanish.
The terminal graph determines the two counts,
so reordering the same graph's edge insertions
leaves the net values unchanged.
\end{proof}

\section{Concrete causal tests and acceptance of the spatial experiment}
For a connected $3$-vertex path, $(|V|,|E|,c)=(3,2,1)$:
the two smooth edge classes are simultaneously Hodge
only for $\tau_1=\tau_2=\tau_3$;
the singular graph realization has $W_0H^1=0$.
For a connected triangle, $(3,3,1)$:
the same two independent Hodge-alignment conditions
remain, while the topology-folded nodal realization
births exactly one weight-zero direction.
For the complete $5$-vertex graph, $(5,10,1)$:
the smooth Hodge-alignment locus has complex
codimension $4$; the nodal topology fold births
$6$ weight-zero directions. Both sum to the
actual ten causal edges.

The corresponding Wolfram computation enumerates
every labelled simple graph on $1\le |V|\le5$
($1+2+8+64+1024=1099$ graphs) and independently
checks $\mathrm{rank} B_\Gamma=|V|-c$,
$\dim\operatorname{coker}B_\Gamma=|E|-|V|+c$,
the constant-modulus alignment rule, and the
exterior-product period
$\tau_w-\tau_v$ with exact arithmetic. It is a
finite consistency check of equations derived above.

\section{Direct connection to the stated all-original-$X$ research goal}
This experiment has now \emph{both} outcomes the
owner requested: marked rational Hodge classes
can lose their Hodge type during a smooth-projective
complex-geometric deformation while the entire
pure Hodge structure survives; and actual singular
causal topology folds birth weight-zero classes
which prevent the result from carrying a pure
weight-one Hodge structure. No Hodge conjecture
hypothesis is used at any point.

A \emph{positive} geometry-to-algebraicity theorem
is also obtained, without any rational Hodge-to-Chow
selector: for each edge class satisfying its exact
Hodge-alignment condition, a genuine algebraic
diagonal realizes its canonical Betti class.
This is a precise theorem for a geometrically
generated edge family, stable under product with
every arbitrary smooth projective $X$.

The proposed implication that \emph{the entire pure
Hodge structure of every smooth projective
realization automatically algebraizes every
rational diagonal Hodge class of arbitrary $X$}
is a logically stronger all-$X$ statement.
Neither the experimental period equation (4)
nor the causal complementarity identity (12)
asserts that stronger implication. In fact
pure Hodge structure is present in both the
aligned and misaligned smooth-projective
realizations, so purity cannot distinguish the
two outcomes \emph{without the additional
geometric alignment law}. Establishing a general
GST geometric alignment-to-algebraic realization
law beyond this edge family would be new
substantive mathematics; we do not assume it.

This keeps the promised focus on whole-space
topology, emergent dimension, Hodge birth/break
and exact geometric realization while clearly
separating the proved geometric implications
from universal conjectural consequences.
\end{document}
